% ===========================================================================
\chapter{Wrong turns}\label{ch:wrongturn}

Every driver misses a turn now and then. A navigator can't stop them; it
can only notice, quickly, and work out a new way from wherever they have
ended up. ``Recalculating\ldots'' This chapter teaches the navigator to do
that, and then builds a pretend traveller, with a pretend GPS, to try it
out on.

\section{Where is the traveller now?}

When the traveller has left the route, the navigator's progress along it
means nothing any more. All it has is the last GPS fix. So, as before, it
assumes the traveller is on a road, but this time any road of the map, not
just the route: it finds the edge nearest to the fix.

Each two-way stretch of road is two edges, one each way, at exactly the
same distance from the fix. Which way is the traveller going? Away from
the route, most likely, since that is how they come to be off it. So of the
two, we take the edge whose end is further from the point where the
traveller left the route.

\begin{append}{navigation.nex}
-- The edge of g nearest p, going away from `behind`: of the two
-- directions along the nearest stretch of road, the one whose end is
-- further from `behind`.
function rejoin(g: Road_Graph, p: Point, behind: Point): Edge
require
  has_edges: g.edge_count > 0
do
  let all: Array[Edge] := []
  across g.nodes as id do
    all := all.concat(g.edges_from(id))
  end
  result := all.get(0)
  let best: Real := distance_to_segment(p, g.island.point(result.source),
    g.island.point(result.target))
  across all as e do
    let a: Point := g.island.point(e.source)
    let b: Point := g.island.point(e.target)
    let d: Real := distance_to_segment(p, a, b)
    let better: Boolean := d < best - 0.5
    if not better and (d - best).abs <= 0.5 then
      -- The same stretch of road the other way: keep the direction that
      -- leads away from `behind`.
      better := b.distance_to(behind) > g.island.point(result.target).distance_to(behind)
    end
    if better then
      best := d
      result := e
    end
  end
end
\end{append}

Two distances within half a metre of each other count as a tie, since two
edges along the same stretch of road should be exactly equal, but rounding
might make them differ in the last decimal place.

\section{Planning again}

The new route starts with the edge the traveller is on, and continues with
the best way from the end of that edge to the destination, planned with
the same algorithm as before. The traveller is already part of the way
along the first edge, so the progress starts there: the distance from the
edge's start to the point of the edge nearest the fix.

\begin{addbefore}{navigation.nex}{to Navigator, before is_off_route}{  is_off_route(): Boolean do}
  -- Plan a new route from where the traveller is now.
  --
  -- The navigator knows only the GPS fixes. The traveller is on the road
  -- nearest the last fix, heading away from where they left the route. So
  -- the new route starts with that stretch of road, and continues with
  -- the best way from the end of it to the destination.
  reroute() do
    let left_at: Point := route.point_at(progress)
    let e: Edge := rejoin(graph, last_fix, left_at)
    let goal: Integer := route.nodes.get(route.nodes.length - 1)
    let rest: Search_Result := find_route(graph, e.target, goal, algorithm, false)
    if rest.found() then
      let path: Array[Integer] := [e.source].concat(rest.path)
      route := create Route_Line.make(graph, path)
      instructions := directions(graph, path, destination)
      let a: Point := graph.island.point(e.source)
      progress := a.distance_to(closest_point_on_segment(last_fix, a,
        graph.island.point(e.target)))
      misses := 0
      reroutes := reroutes + 1
    end
  end
\end{addbefore}

The new route may well begin by turning round: if the traveller has just
gone the wrong way down a long road with no turnings, the quickest way back
may be a U-turn. The directions will say so: ``Turn sharp left''.

\section{A pretend traveller}

To watch all this happen, we need someone to navigate. Our traveller
moves along a path of nodes, at the speed of each road: an edge of 100
metres that takes 12 seconds to drive is covered at $100/12$ metres per
second. The traveller knows where it really is: on the edge from
\code{path[leg]} to \code{path[leg + 1]}, \code{offset} metres along it.

\begin{newfile}{simulation.nex}
-- simulation.nex: a pretend traveller with a pretend GPS, so that the
-- navigator has someone to guide.

intern Navigation
intern Rng

-- The traveller, moving along a path of nodes at the speed of each road.
-- Usually the path is the navigator's route, but a traveller can take a
-- wrong turn.
class Traveller
create
  make(g: Road_Graph, nodes: Array[Integer])
  require
    is_a_route: is_route(g, nodes)
    can_move: nodes.length >= 2
  do
    graph := g
    path := nodes
    leg := 0
    offset := 0.0
    wrong_turn_pending := false
  end
feature
  graph: Road_Graph
  -- The nodes the traveller means to pass through.
  path: Array[Integer]
  -- The traveller is `offset` metres along the edge from path[leg] to
  -- path[leg + 1].
  leg: Integer
  offset: Real
  -- Set to make the traveller go the wrong way at the next junction.
  wrong_turn_pending: Boolean

  has_arrived(): Boolean do
    result := leg >= path.length - 1
  end

  -- The edge the traveller is on.
  edge(): Edge
  require
    on_the_way: not has_arrived()
  do
    if ?graph.edge_between(path.get(leg), path.get(leg + 1)) as e then
      result := e
    else
      raise "the traveller's path has a step that is not an edge"
    end
  end

  -- Where the traveller really is.
  position(): Point do
    if has_arrived() then
      result := graph.island.point(path.get(path.length - 1))
    else
      let a: Point := graph.island.point(path.get(leg))
      let b: Point := graph.island.point(path.get(leg + 1))
      result := a + (b - a) * (offset / a.distance_to(b))
    end
  end

  -- Which way the traveller is facing, as a compass bearing.
  heading(): Real do
    if path.length < 2 then
      result := 0.0
    else
      let i: Integer := leg.min(path.length - 2)
      result := bearing(graph.island.point(path.get(i)), graph.island.point(path.get(i + 1)))
    end
  end

  -- Go the wrong way at the next junction.
  take_wrong_turn() do
    wrong_turn_pending := true
  end
\end{newfile}

Moving on by some number of seconds means covering the rest of the current
edge, if there is time, then the next edge, and so on, until the time is
used up or the traveller arrives. Each time the traveller reaches a node,
if a wrong turn is due, it takes it there.

\begin{continued}{simulation.nex}
  -- Move on by `seconds` of travel.
  advance(seconds: Real)
  require
    non_negative: seconds >= 0.0
  do
    let left: Real := seconds
    from
    until
      left <= 0.0 or has_arrived()
    do
      let e: Edge := edge()
      let speed: Real := e.length / e.seconds
      let to_go: Real := (e.length - offset) / speed
      if left < to_go then
        offset := offset + left * speed
        left := 0.0
      else
        left := left - to_go
        offset := 0.0
        leg := leg + 1
        if wrong_turn_pending and not has_arrived() then
          go_wrong()
        end
      end
    end
  end
\end{continued}

When the navigator plans a new route, the traveller takes it. The new route
starts with the edge the navigator believes the traveller is on; if that is
the edge the traveller really is on, the traveller simply swaps its path
for the new one. If the traveller has already reached the end of that edge,
the new route starts where it is about to arrive, which works too. If
neither, the navigator has guessed wrong, and the traveller carries on
regardless. The navigator will soon find it off the new route as well, and
try again.

\begin{continued}{simulation.nex}
  -- Take the navigator's new route, which starts on the traveller's edge
  -- or at the node the traveller is about to reach.
  follow(nodes: Array[Integer]) do
    if has_arrived() then
      if nodes.get(0) = path.get(path.length - 1) and nodes.length >= 2 then
        path := nodes
        leg := 0
        offset := 0.0
      end
    elseif nodes.length >= 2 and nodes.get(0) = path.get(leg)
        and nodes.get(1) = path.get(leg + 1) then
      path := nodes
      leg := 0
    elseif nodes.get(0) = path.get(leg + 1) then
      path := [path.get(leg)].concat(nodes)
      leg := 0
    end
  end
\end{continued}

A wrong turn happens at a junction: a node with some way out other than
the planned one and the one the traveller came from. Most nodes of a route
are not junctions, just points along a road, so a wrong turn that is due
waits for the next real junction. There, the traveller takes whichever
other road turns least, the classic mistake of going straight on when the
route turns, and then simply keeps going as straight as it can, for a
dozen edges, wherever that leads.

\begin{continued}{simulation.nex}
private feature
  -- At path[leg], turn off the planned way onto another road, and keep
  -- going straight on, wherever that leads.
  go_wrong() do
    let here: Integer := path.get(leg)
    let came_from: Integer := path.get(leg - 1)
    let planned: Integer := path.get(leg + 1)
    let options: Array[Edge] := []
    across graph.edges_from(here) as e do
      if e.target /= came_from and e.target /= planned then
        options.add(e)
      end
    end
    if not options.is_empty then
      let wrong: Edge := straightest(came_from, here, options)
      let detour: Array[Integer] := [wrong.target]
      let previous: Integer := here
      repeat 12 do
        let at: Integer := detour.get(detour.length - 1)
        let onward: Array[Edge] := []
        across graph.edges_from(at) as e do
          if e.target /= previous then
            onward.add(e)
          end
        end
        if not onward.is_empty then
          detour.add(straightest(previous, at, onward).target)
          previous := at
        end
      end
      path := path.take(leg + 1).concat(detour)
      wrong_turn_pending := false
    end
  end

  -- Of `options`, the edges leaving `at`, the one that turns least for a
  -- traveller arriving from `before`.
  straightest(before, at: Integer, options: Array[Edge]): Edge
  require
    has_options: not options.is_empty
  do
    let island: Island_Map := graph.island
    let heading_in: Real := bearing(island.point(before), island.point(at))
    result := options.get(0)
    across options as e do
      let turn: Real := turn_angle(heading_in, bearing(island.point(at), island.point(e.target))).abs
      let best: Real := turn_angle(heading_in, bearing(island.point(at),
        island.point(result.target))).abs
      if turn < best then
        result := e
      end
    end
  end
end
\end{continued}

\section{A pretend GPS}

A real GPS fix is off by a few metres, in a random direction, by a random
amount: usually a little, occasionally more. Errors like that follow the
\emph{normal distribution}, the bell curve. To make them, we need random
numbers, and to be able to repeat a run exactly (so that the pictures in
this tutorial come out the same every time), the random numbers must come
from a generator we can \emph{seed}: start from a chosen number, and get
the same sequence every time.

\begin{newfile}{rng.nex}
-- rng.nex: a small random number generator that can be seeded, for the
-- simulated GPS. With a seed, the "random" numbers are the same on every
-- run, so a simulated trip can be repeated exactly.
--
-- This is the Park-Miller "minimal standard" generator:
--     state := state * 48271 mod (2^31 - 1)

class Rng
create
  make(seed: Integer) do
    state := 1 + seed.abs % (MODULUS - 1)
  end
feature
  MODULUS = 2147483647
  MULTIPLIER = 48271

  state: Integer

  -- A real number spread evenly over the open interval (0, 1).
  uniform(): Real
  do
    state := (state * MULTIPLIER) % MODULUS
    result := state.to_real / MODULUS.to_real
  ensure
    open_interval: result > 0.0 and result < 1.0
  end

  -- A number from the standard normal ("bell curve") distribution, with
  -- mean 0 and standard deviation 1, by the Box-Muller transform.
  gauss(): Real do
    let radius: Real := (-2.0 * uniform().log).sqrt
    let angle: Real := 2.0 * Real.PI * uniform()
    result := radius * angle.cos
  end
invariant
  state_in_range: state >= 1 and state < MODULUS
end
\end{newfile}

The generator multiplies its state by 48,271 and keeps the remainder after
dividing by $2^{31} - 1$, a prime number. The result jumps about so
thoroughly that it looks random, and dividing by the modulus gives a number
between 0 and 1.

The \emph{Box--Muller transform} turns two such evenly spread numbers,
$u_1$ and $u_2$, into one from the bell curve: $\sqrt{-2\ln u_1}
\cos(2\pi u_2)$. It is a neat piece of mathematics: picking a random angle,
$2\pi u_2$, and a random distance, $\sqrt{-2\ln u_1}$, gives a point whose
$x$ and $y$ both follow the bell curve. We use just the $x$.

A fix is then the true position plus a normal error in each direction,
four metres on average:

\begin{append}{simulation.nex}
-- A pretend GPS receiver. Every fix is the true position plus a random
-- error of a few metres, the way a real one is.
class Gps
create
  make(seed: Integer, metres: Real) do
    rng := create Rng.make(seed)
    error := metres
  end
feature
  rng: Rng
  -- The typical size of the error, in metres (its standard deviation).
  error: Real

  fix(truth: Point): Point do
    result := truth + create Point.make(rng.gauss() * error, rng.gauss() * error)
  end
end
\end{append}

\section{A trip}

Finally, a trip puts the traveller, the GPS and the navigator together, with
a clock. As the clock moves on, the traveller moves, and once a second the
GPS takes a fix and gives it to the navigator. If the navigator finds the
traveller off the route, it plans a new one, and the traveller takes it.

Notice how little the navigator is told. It never asks the traveller where
it is: it gets only fixes, with their errors, exactly as a real navigator
gets only what the GPS chip says.

\begin{append}{simulation.nex}
-- A trip: a traveller, their GPS, and the navigator guiding them, with
-- the clock that moves them all on together.
class Trip
create
  make(g: Road_Graph, path: Array[Integer], destination, algorithm: String, seed: Integer)
  require
    is_a_route: is_route(g, path)
    can_move: path.length >= 2
  do
    traveller := create Traveller.make(g, path)
    gps := create Gps.make(seed, 4.0)
    navigator := create Navigator.make(g, path, destination, algorithm)
    clock := 0.0
    next_fix := 0.0
    fix := traveller.position()
  end
feature
  traveller: Traveller
  gps: Gps
  navigator: Navigator
  -- Seconds since the trip began...
  clock: Real
  -- ...and when the GPS gives its next fix (once a second).
  next_fix: Real
  -- The latest fix.
  fix: Point

  -- Move the trip on by `seconds`.
  advance(seconds: Real) do
    let stop: Real := clock + seconds
    from
    until
      next_fix > stop
    do
      traveller.advance(next_fix - clock)
      clock := next_fix
      fix := gps.fix(traveller.position())
      navigator.update(fix)
      if navigator.is_off_route() then
        navigator.reroute()
        traveller.follow(navigator.route.nodes)
      end
      next_fix := next_fix + 1.0
    end
    traveller.advance(stop - clock)
    clock := stop
  end

  is_over(): Boolean do
    result := traveller.has_arrived() and navigator.has_arrived()
  end
end
\end{append}

\code{advance} can be asked to move on by any amount of time, a fraction of
a second or a minute, and it still takes a fix at every whole second in
between. The live map will move it on by whatever time has passed since
the last frame it drew.

\section{Taking a wrong turn}

This script drives from the harbour to the lighthouse. Forty seconds in,
it tells the traveller to take a wrong turn, and it prints what the
navigator says along the way: the prompt whenever the next instruction
changes, and every thirty seconds in between.

\begin{tryit}{try_trip.nex}
intern Simulation

let island: Island_Map := create Island_Map.load(user_path("kestrel_island.map"))
let g: Road_Graph := create Road_Graph.build(island, Mode.Driving)
if ?island.find_place("harbour") as a and ?island.find_place("lighthouse") as b then
  let found: Search_Result := find_route(g, g.nearest_node(island.point(a.node_id)),
    g.nearest_node(island.point(b.node_id)), "astar", false)
  let trip: Trip := create Trip.make(g, found.path, b.name, "astar", 7)
  console.print_line("   0 s  " + trip.navigator.prompt())
  let step: String := trip.navigator.next_instruction().text
  let told: Boolean := false
  from
  until
    trip.is_over() or trip.clock > 1000.0
  do
    if trip.clock >= 40.0 and not told then
      trip.traveller.take_wrong_turn()
      told := true
    end
    let before: Integer := trip.navigator.reroutes
    trip.advance(1.0)
    let time: String := trip.clock.format(0).pad_start(" ", 4) + " s  "
    if trip.navigator.reroutes > before then
      console.print_line(time + "Off the route: planning a new one")
    end
    let next: String := trip.navigator.next_instruction().text
    if next /= step or trip.clock.round % 30 = 0 then
      console.print_line(time + trip.navigator.prompt().pad_end(" ", 44)
        + distance_text(trip.navigator.remaining_distance()) + " to go")
      step := next
    end
  end
  console.print_line(trip.clock.format(0).pad_start(" ", 4) + " s  " + trip.navigator.prompt())
  console.print_line("The trip took " + trip.clock.format(0) + " seconds; the route planned at the start takes "
    + (create Route_Line.make(g, found.path).duration()).format(0) + ".")
end
\end{tryit}
\tryoutput{try_trip}

Read it from the top. The traveller sets off up Harbour Street, and at 40
seconds a wrong turn is due. At the next junction the traveller turns off
the planned route. A few seconds later the navigator, seeing only the fixes,
has noticed: three fixes in a row too far from the route. It finds the road
the traveller is now on, plans a new route from there, and starts giving
instructions for it straight away: back up Net Lane to the Coast Road, and
on to Hill Road. The detour costs about half a minute, and the traveller
still gets to the lighthouse.

\section{Checking the navigation}

The last checks follow a whole trip, twice: once as planned, which must
arrive without ever leaving the route, in about the time the route should
take; and once with a wrong turn, which must still arrive, with at least
one new route planned on the way. They also check the route line: that it
starts and ends at the right nodes, and that a point 10 metres to the side
of it is matched back onto it, at the right place.

\begin{addbefore}{checks.nex}{before main}{function main() do}
function check_navigation(c: Checker, m: Island_Map, g: Road_Graph) do
  if ?m.find_place("harbour") as a and ?m.find_place("lighthouse") as b then
    let found: Search_Result := find_route(g, node_of(g, a), node_of(g, b), "astar", false)
    let line: Route_Line := create Route_Line.make(g, found.path)
    c.close_to("a route line starts at the start", 0.0,
      line.point_at(0.0).distance_to(m.point(found.path.get(0))), 0.000001)
    c.close_to("and ends at the end", 0.0, line.point_at(line.length())
      .distance_to(m.point(found.path.get(found.path.length - 1))), 0.000001)
    c.close_to("the time to the end is the duration", line.duration(),
      line.time_at(line.length()), 0.000001)
    let halfway: Real := line.length() / 2.0
    let beside: Point := line.point_at(halfway)
      + direction(line.heading_at(halfway) + 90.0) * 10.0
    let found_again: Line_Match := line.locate(beside, 0.0, line.length())
    c.close_to("a point 10 m off the route is matched to it", 10.0, found_again.offset, 0.001)
    c.close_to("at the right place along it", halfway, found_again.along, 0.001)

    let trip: Trip := create Trip.make(g, found.path, b.name, "astar", 1)
    from
    until
      trip.is_over() or trip.clock > 3.0 * line.duration()
    do
      trip.advance(1.0)
    end
    c.check("a simulated trip arrives", trip.is_over())
    c.check("without leaving the route", trip.navigator.reroutes = 0)
    c.close_to("in the time the route should take", line.duration(), trip.clock, 3.0)

    let lost: Trip := create Trip.make(g, found.path, b.name, "astar", 2)
    lost.traveller.take_wrong_turn()
    from
    until
      lost.is_over() or lost.clock > 3.0 * line.duration()
    do
      lost.advance(1.0)
    end
    c.check("after a wrong turn the traveller still arrives", lost.is_over())
    c.check("because the navigator planned a new route", lost.navigator.reroutes >= 1)
  end
end
\end{addbefore}

\begin{tightaddbefore}{checks.nex}{to main, before c.summary()}{  c.summary()}
  check_navigation(c, m, driving)
\end{tightaddbefore}

\begin{tightaddbefore}{checks.nex}{at the top}{intern Directions}
intern Simulation
\end{tightaddbefore}

\runoutput{checks.nex}{checks}

All 56 checks pass.
